\section{Complete paths and their selection}
\label{sec:main-results}

We stop the original trajectory at its first crossing of one fixed
downstream section. The limiting object retains the incoming rigid
arc, every return in an accumulating cascade and the final contact
or flight prefix. Its contact pieces carry an order, although their
physical durations vanish.

\subsection{The curves and their topology}
For continuous oriented curves $f,g:[0,1]\to\R^5$, put
\begin{equation}
 \dF(f,g)=\inf_{h\in\mathrm{Homeo}_+([0,1])}
                 \sup_{q\in[0,1]}|f(h(q))-g(q)|.
 \label{eq:ps5-frechet}
\end{equation}
We identify curves at zero distance. Pauses and constant pieces are
therefore allowed. Equivalently, one can use two nondecreasing
continuous surjections in the comparison: approximate a surjection
$a$ by $(1-\eta)a+\eta\,\mathrm{id}$ and use uniform continuity,
then compose the resulting homeomorphisms. For compact path images,
\begin{equation}
 \dH(A,B)=\max\{\sup_{x\in A}\inf_{y\in B}|x-y|,
                    \sup_{y\in B}\inf_{x\in A}|x-y|\}
 \le\dF(f,g),\quad A=f([0,1]),\ B=g([0,1]).
 \label{eq:ps5-hausdorff}
\end{equation}

Write $I_{\lambda,\nu}$ for the prescribed saturated rigid incoming
arc from $(\phi_\nu,0,\omega_\nu,0,0)$ to $G$, oriented in physical
time. Appendix~\ref{sec:detuned-return} and Section~\ref{sec:graph} construct an
exact limiting graph return $F=F_\lambda$ and its ordered one-return
curve $\cC_{\lambda,d}:[0,3]\to\R^5$. The three parameter intervals
describe saturated compression, inherited normal release and the
following first free flight: $\alpha=0,1,2,3$ mark touchdown, slip
arrest, release and the next touchdown, respectively.
These are finite-scale curves, retaining
their exact $\sqrt d$ dependence. They satisfy
\begin{equation}
 \cC_{\lambda,d}(3)=\cC_{\lambda,F(d)}(0),\qquad
 \sup_{0\le\alpha\le3}|\cC_{\lambda,d}(\alpha)-G|
                      \le C d^{3/4}.
 \label{eq:ps5-joins}
\end{equation}
On contact $p(\cC_{\lambda,d}(\alpha))=-ad$; on the free part it
strictly decreases from $-ad$ to $-aF(d)$. For fixed sufficiently
small $s>0$, set $d_*=F^{-1}(s)$. There is consequently a unique
continuous free-section parameter $j_s:[d_*,s]\to[2,3]$ with
\begin{equation}
 p(\cC_{\lambda,d}(j_s(d)))=-as,\qquad
 j_s(d_*)=3,\quad j_s(s)=2.
 \label{eq:path-free-section}
\end{equation}
The possible terminal data form the compact set
\begin{equation}
 \mathcal A_s=\{(d,j_s(d)):d_*\le d\le s\}
                       \cup(\{s\}\times[0,2]).
 \label{eq:ps5-domain}
\end{equation}
Its graph branch records free-flight stopping; the vertical segment
records contact stopping.

For $d\in[d_*,s]$, let $b_m=F^{-m}(d)$, $m\ge0$. These bases tend
to zero. Define $L_d(0)=G$ and trace $\cC_{\lambda,b_m}$ linearly on
$[2^{-m},2^{-(m-1)}]$ for $m\ge1$. The joins and bound in
\eqref{eq:ps5-joins} give a continuous curve. The complete curves are
\begin{equation}
 \cW_{\lambda,\nu,s}(d,\alpha)
   =I_{\lambda,\nu}\star L_d\star
                   \cC_{\lambda,d}|_{[0,\alpha]},
       \qquad(d,\alpha)\in\mathcal A_s.
 \label{eq:ps5-whole}
\end{equation}
Here $\star$ preserves order and joins equal endpoints. One may assign
one third of $[0,1]$ to each factor; a zero final prefix is constant.
Its compact image is
\begin{equation}
 W_{\lambda,\nu,s}(d,\alpha)
  =I_{\lambda,\nu}([0,1])\cup\{G\}
    \cup\bigcup_{m\ge1}\cC_{\lambda,b_m}([0,3])
    \cup\cC_{\lambda,d}([0,\alpha]).
 \label{eq:ps5-whole-image}
\end{equation}
The parameter $\alpha$ records the order along each completed contact
and free return.
At admission the two terminal pairs represent one ordered path:
$\cW_{\lambda,\nu,s}(d_*,3)=\cW_{\lambda,\nu,s}(s,0)$.
Indeed, $F(d_*)=s$ identifies the final full return of the former
with the last return already contained in $L_s$. Their parameter
pairs remain distinct. Figure~\ref{fig:ordered-stopped-path} shows
the order and terminal domain.
% MIXED: exact labels and joins, schematic widths and free-graph shape.
% The upper strip is an order diagram, not a coordinate projection or clock.
\begin{figure}[!htbp]
\centering
\begin{tikzpicture}[
 x=1cm,y=1cm,font=\small,
 every node/.style={inner sep=1.5pt},
 box/.style={draw=black!70,line width=.5pt,fill=black!3},
 arrow/.style={-{Latex[length=1.6mm,width=1.1mm]},line width=.5pt},
 guide/.style={draw=black!35,densely dotted,line width=.45pt}
]
 \node[anchor=west,font=\small\bfseries] at (0,5.75)
   {(a) The complete ordered path};
 \draw[arrow] (0,5.30) -- (14.4,5.30)
   node[midway,above] {order along the limiting curve};
 \draw[box] (0,4.45) rectangle (2.05,5.00);
 \node at (1.025,4.725) {$I_{\lambda,\nu}$};
 \fill (2.05,4.725) circle[radius=1.3pt];
 \node[above right] at (2.05,4.95) {$G$};
 \node at (2.60,4.725) {$\cdots$};
 \draw[box] (3.05,4.45) rectangle (4.60,5.00);
 \draw[box] (4.60,4.45) rectangle (6.15,5.00);
 \draw[box] (6.15,4.45) rectangle (7.70,5.00);
 \node at (3.825,4.725) {$\mathcal C_{\lambda,b_3}$};
 \node at (5.375,4.725) {$\mathcal C_{\lambda,b_2}$};
 \node at (6.925,4.725) {$\mathcal C_{\lambda,b_1}$};
 \node at (8.15,4.725) {$\star$};
 \draw[box] (8.60,4.45) rectangle (14.40,5.00);
 \node at (11.50,4.725)
   {$\mathcal C_{\lambda,d}|_{[0,\alpha]}$ \ (final prefix)};
 \node at (5.15,4.15) {$L_d:\quad b_m=F_\lambda^{-m}(d)\longrightarrow0$};
 \node[anchor=west] at (0,3.62) {One return:};
 \draw[box,fill=black!9] (2.05,3.28) rectangle (5.85,3.88);
 \draw[box] (5.85,3.28) rectangle (10.40,3.88);
 \draw[box,fill=white] (10.40,3.28) rectangle (14.40,3.88);
 \node at (3.95,3.58) {compression $[0,1]$};
 \node at (8.125,3.58) {inherited release $[1,2]$};
 \node at (12.40,3.58) {free flight $[2,3]$};

 \node[anchor=west,font=\small\bfseries] at (0,2.77)
   {(b) The terminal domain $\mathcal A_s$};
 \draw[arrow] (.50,.05) -- (6.65,.05) node[right] {$d$};
 \draw[arrow] (.50,.05) -- (.50,2.30) node[above] {$\alpha$};
 \foreach \yy/\ll in {.05/0,.70/1,1.35/2,2.00/3}{
   \draw (.44,\yy) -- (.56,\yy);
   \node[left=2pt] at (.44,\yy) {$\ll$};
 }
 \draw[guide] (1.50,.05) -- (1.50,2.00) -- (.50,2.00);
 \draw[guide] (.50,1.35) -- (5.45,1.35);
 \node[below=3pt] at (1.50,.05) {$d_*$};
 \node[below=3pt] at (5.45,.05) {$s$};
 % Only the continuous graph and its endpoint values are asserted.
 \draw[densely dashed,line width=.8pt]
   (1.50,2.00) .. controls (2.80,1.98) and (4.20,1.40) .. (5.45,1.35);
 \node[above] at (3.50,1.85) {$(d,j_s(d))$};
 \draw[line width=.8pt] (5.45,.05) -- (5.45,1.35);
 \node[rotate=90,above=3pt] at (5.45,.70) {contact};
 \fill (1.42,1.92) rectangle (1.58,2.08);
 \fill (5.37,-.03) rectangle (5.53,.13);
 \draw[fill=white,line width=.7pt] (5.45,1.35) circle[radius=1.8pt];
 \fill (5.45,.70) circle[radius=1.5pt];
 \node[above=3pt] at (1.50,2.00) {$(d_*,3)$};
 \node[right=3pt] at (5.45,1.35) {$(s,2)$};
 \node[right=3pt] at (5.45,.70) {$(s,1)$};
 \node[above right=3pt] at (5.45,.05) {$(s,0)$};
 \node[draw=black!50,rounded corners=2pt,align=center,
       text width=6.0cm,inner sep=6pt] at (11.10,1.45)
   {Two parameter pairs, one ordered path:\\[3pt]
    $\mathcal W_{\lambda,\nu,s}(d_*,3)
       =\mathcal W_{\lambda,\nu,s}(s,0)$\\[3pt]
    $d_*=F_\lambda^{-1}(s)$};
\end{tikzpicture}
\caption{Complete stopped paths and their terminal data.
(a) The incoming arc is followed by returns accumulating at $G$
and a final prefix; the displayed $b_3,b_2,b_1$ are the last three
full-return bases. Strip widths encode order, not duration.
(b) The free graph and contact segment form $\mathcal A_s$; the
graph shape is schematic. Arrest is $(s,1)$. Matching squares
represent the single endpoint identification of complete ordered
paths, while the parameter pairs remain distinct. Identifying the
two interval endpoints gives the quotient circle of
Corollary~\ref{cor:accumulation-circle}.}
\label{fig:ordered-stopped-path}
\label{fig:terminal-parameter-domain}
\end{figure}



\subsection{Classification, common selection and contact sensitivity}
Let $t_n$ be the successive actual touchdowns after the critical
release and put $\delta_n=-p(t_n)/a$ whenever $p(t_n)<0$. The second
and subsequent touchdowns have this sign. Let $t_s$ denote the
actual first crossing of $p=-as$. Existence of the required events
is part of the following theorem. For sufficiently small compliance,
put
\begin{equation}
 \begin{gathered}
 N_s=\min\{n\ge3:\delta_n\ge s\},\qquad
 d_\eps=\delta_{N_s-1},\qquad
 \alpha_\eps=h_\eps^{-1}(t_s),\\
 f_\eps(q)=\cX_\eps(t_s q),\qquad H_\eps=f_\eps([0,1]).
 \end{gathered}
 \label{eq:ps5-stop}
\end{equation}
The increasing pairing $h_\eps:[0,3]\to[t_{N_s-1},t_{N_s}]$
is constructed in Lemma~\ref{path:ordered-return}. Its compression
parameter is the fraction of actual normal impulse up to the first
$u_r=M_{\rm sl}\eps$ event; the release parameter is the elapsed fraction
from that event to the first force zero. The last interval uses
elapsed free time. Thus $\alpha_\eps$ is read from the actual return
containing the first section time.

\begingroup
\renewcommand{\thetheorem}{A}
\begin{theorem}[Classification and actual realization]
\label{thm:whole}
Under the hypotheses of Section~\ref{sec:model}, there is a common
$s_0>0$ such that, for every fixed $0<s\le s_0$ and all sufficiently
small $\eps$, the original solution reaches its first section time
$t_s$ transversely, and $N_s\to\infty$ as $\eps\downarrow0$.
For each fixed
complete label $\mathfrak l=(\lambda,k,\theta,\nu,\mathbf n)$, the
family $f_\eps$ is relatively compact for \eqref{eq:ps5-frechet},
and
\[
 \Clust_{\eps\downarrow0}(d_\eps,\alpha_\eps)=\mathcal A_s.
\]
Its ordered-path accumulation set is exactly
\[
 \{\cW_{\lambda,\nu,s}(d,\alpha):(d,\alpha)\in\mathcal A_s\}.
\]
Every listed curve is attained by an actual sequence of compliances
with this label fixed. More precisely, for every sequence $\eps_i\downarrow0$ with
$(d_{\eps_i},\alpha_{\eps_i})\to(d,\alpha)$, one has
$(d,\alpha)\in\mathcal A_s$ and
\begin{equation}
 \dF(f_{\eps_i},\cW_{\lambda,\nu,s}(d,\alpha))\to0,\qquad
 \dH(H_{\eps_i},W_{\lambda,\nu,s}(d,\alpha))\to0.
 \label{eq:ps5-whole-limit}
\end{equation}
The Hausdorff accumulation set is exactly the corresponding image
family in \eqref{eq:ps5-whole-image}.
\end{theorem}
\endgroup

For fixed $\lambda,\nu,s$, this accumulation family is independent
of the profile, belt coefficient and fast preparation within the
prescribed classes. Its contact segment records every possible
stopping point of the terminal contact. A specified free
target has the stronger selection property below.

\begingroup
\renewcommand{\thetheorem}{B}
\begin{theorem}[Common selection of a free target]
\label{thm:common-free}
With the compact classes and $s_0$ of Theorem~\ref{thm:whole}, fix
$\lambda$, $0<s\le s_0$ and $d\in(d_*,s)$. The normalized graph phase
$\mathcal S_\lambda$ gives the explicit sequence
\begin{equation}
 \eps_m(d)=\exp\{-\tfrac{32}{9}(3/4)^{-m}
                                      \mathcal S_\lambda(d)\}.
 \label{eq:ps5-common}
\end{equation}
Along it,
\begin{equation}
 \sup_{(k,\theta,\nu,\mathbf n)}
 \left\{|(d_{\eps_m},\alpha_{\eps_m})-(d,j_s(d))|
 +\dF(f_{\eps_m},\cW_{\lambda,\nu,s}(d,j_s(d)))\right\}\to0,
 \label{eq:common-target-limit}
\end{equation}
where the supremum is over the prescribed compact auxiliary class.
The corresponding Hausdorff image convergence is uniform as well.
The two free endpoints $(d_*,3)$ and $(s,2)$ have the same
uniform convergence along the explicit sequences in
Proposition~\ref{sel:threshold-selection} below.
\end{theorem}
\endgroup

The sequence fixes the target while the auxiliary labels vary;
each incoming datum retains its own arc $I_{\lambda,\nu}$.
The exact phase $\mathcal S_\lambda$ contains the finite-scale
return information needed at the fixed section $s$.
In particular, the explicit sequence uses this constructed phase;
replacing it by $\log(K_\lambda^4/d)$ at fixed $d$ need not select
the same target.

\begingroup
\renewcommand{\thetheorem}{C}
\begin{theorem}[Contact selection by the entrance constant]
\label{thm:contact-selection}
Under the hypotheses of Section~\ref{sec:model}, fix $\lambda$ and
$0<s\le s_0$ sufficiently small before taking the compliance limit.
The original physical preparation determines a continuous bounded
entrance constant $\chi_\lambda(\mathfrak a)$ for the prescribed
auxiliary labels $\mathfrak a=(k,\theta,\nu,\mathbf n)$; its
construction is given in Proposition~\ref{loc:entrance} and
\eqref{sel:chi}.
\begin{enumerate}[label=\textup{(\roman*)},leftmargin=*]
\item For every two fixed auxiliary labels
$\mathfrak a_1,\mathfrak a_2$ with unequal entrance constants, there
is $\eps_0>0$, depending on this pair and the fixed data, such that
for $0<\eps<\eps_0$ their first section times cannot both satisfy
\[
 t_{j_i}^{\mathfrak a_i}<t_s^{\mathfrak a_i}
       \le t_{L,j_i}^{\mathfrak a_i},\qquad i=1,2,
\]
where $t_{L,j_i}^{\mathfrak a_i}$ is the first release after
touchdown $j_i$. Consequently no common compliance sequence selects
contact-interior terminal pairs $(s,\alpha_i)$ for both labels,
where $0<\alpha_i<2$, or the corresponding ordered paths
$\cW_{\lambda,\nu_i,s}(s,\alpha_i)$, even when
$\alpha_1\ne\alpha_2$.
\item For one fixed auxiliary label $\mathfrak a_0$ and
$0<\alpha_0<2$, suppose that a sequence $\eps_m\downarrow0$ satisfies
\[
 f_{\eps_m}^{\mathfrak a_0}
   \longrightarrow\cW_{\lambda,\nu_0,s}(s,\alpha_0)
   \qquad\hbox{in }\dF.
\]
Along this same given sequence, every other fixed label
$\mathfrak a$ with $\chi_\lambda(\mathfrak a)
<\chi_\lambda(\mathfrak a_0)$ has terminal-pair limit $(d_*,3)$
and ordered-path limit $\cW_{\lambda,\nu,s}(s,0)$; if
$\chi_\lambda(\mathfrak a)>\chi_\lambda(\mathfrak a_0)$, the limits
are $(s,2)$ and $\cW_{\lambda,\nu,s}(s,2)$, respectively.
The corresponding Hausdorff image convergence holds, and each
incoming datum retains its own arc $I_{\lambda,\nu}$.
\end{enumerate}
Equality of entrance constants is necessary for common interior
contact selection; these estimates do not establish sufficiency.
\end{theorem}
\endgroup

\begin{corollary}[Belt-coefficient sensitivity]
\label{cor:belt-contact-selection}
Fix $\lambda,\theta,\nu,\mathbf n$ and a sufficiently small
$0<s\le s_0$ under the hypotheses of Theorem~\ref{thm:whole}.
The entrance constant $\chi_{\lambda,\theta}(k)$ is strictly increasing
in $k$. Thus any two distinct fixed belt coefficients satisfy the
exclusion in Theorem~\ref{thm:contact-selection}\textup{(i)}.
Fix $k_0\in I$ and $0<\alpha<2$. Along any sequence
$\eps_m\downarrow0$ for which
$f_{\eps_m}^{(k_0)}\to\cW_{\lambda,\nu,s}(s,\alpha)$ in $\dF$,
every other fixed $k\in I$ satisfies
\[
\begin{array}{c|c|c}
 &\text{terminal pair limit}&\text{ordered-path limit}\\ \hline
 k<k_0 &(d_*,3)&\cW(d_*,3)=\cW(s,0)\\
 k>k_0 &(s,2)&\cW(s,2).
\end{array}
\]
The corresponding Hausdorff image convergence also holds.
\end{corollary}

Proposition~\ref{sel:threshold-selection}\textup{(i)} gives the
explicit threshold sequences and their touchdown counts in
Section~\ref{sec:contact-selection}; its parts \textup{(ii)} and
\textup{(iii)} prove Theorem~\ref{thm:contact-selection}.
The belt corollary follows
from the strict order of its entrance constant; the same selection
law applies to the locally constructed profile curves in
Corollary~\ref{cor:friction-profile-selection}. Section~\ref{sec:observations}
explains how the observation rule changes common prefix selection.

The complete physical path does not determine every additionally
scaled observable: Corollary~\ref{cor:memory} realizes an interval of
limiting effective-gap velocities $v+z/\eps$ above the same release
endpoint path. Lemma~\ref{lem:terminal-separation} identifies the sole
nontrivial duplication of terminal parameters, and
Corollary~\ref{cor:accumulation-circle} obtains the corresponding
quotient-circle description for terminal states, ordered paths and
compact images, at fixed $\lambda,\nu,s$ in the full physical state.

The proof must start with the original physical preparation and end
with whole curves. We first state the entrance and return estimates,
then obtain one precise comparison from the actual third touchdown.
The realization and selection arguments use that same comparison.
